% Generated by scripts/assemble.py; corrections belong in source/corrections.json.
\examhead{172}{西南交通大学 2026 年数学分析试卷(回忆版)}{163}
% p164-b07; source PDF page 164
\blockanchor{p164-b07}
(注.题目来自考生回忆,部分表述或数据可能有误)\par
% p164-b08; source PDF page 164
\blockanchor{p164-b08}
\question{1}(10 分) 已知 $\{x_{n}\}$ 是无界数列, 但非无穷大量, 证明: 存在 $\{x_{n}\}$ 的两个子列, 其中一个子列是无穷大量, 另一个是收敛子列.\par
% p164-b09; source PDF page 164
\blockanchor{p164-b09}
\question{2}(15 分) 应用 Bolzano-Weierstrass 定理证明闭区间上连续函数的有界性定理.\par
% p164-b10; source PDF page 164
\blockanchor{p164-b10}
\question{3}(10 分) 设 $f_{+}^{\prime}(x_0) > 0$, $f_{-}^{\prime}(x_0) < 0$, 证明: $x_0$ 是 $f(x)$ 的极小值点.\par
% p164-b11; source PDF page 164
\blockanchor{p164-b11}
\question{4}(10 分) 设函数 $f(x) = \begin{cases} x^2, & x \text{ 为有理数;} \\ 0, & x \text{ 为无理数.} \end{cases}$ 证明: $f(x)$ 只在 $x = 0$ 处可导.\par
% p164-b12; source PDF page 164
\blockanchor{p164-b12}
\question{5}(10 分) 讨论反常积分 $\int_{0}^{+\infty} \frac{\ln(1 + x)}{x^n} \, \mathrm{d}x$ 的敛散性, 其中 $n$ 为实数.\par
% p164-b13; source PDF page 164
\blockanchor{p164-b13}
\question{6}(15 分) 讨论下列函数项级数在所指定区间的一致收敛性: $\sum_{n=1}^{\infty} 2^{n} \sin \frac{1}{3^{n} x}, x \in (0, +\infty)$.\par
% p164-b14; source PDF page 164
\blockanchor{p164-b14}
\question{7}(15 分) 判断序列 $f_{n}(x) = \frac{nx}{1 + n^{2}x^{2}} (n = 1,2,\dots)$ 在 $[0,1]$ 上的一致收敛性.\par
% p164-b15; source PDF page 164
\blockanchor{p164-b15}
\question{8}(15 分) 设函数 $f(x, y)$ 在 $x \in [a, A]$, $y \in [b, B]$ 上连续, 函数列 $\{\varphi_n(x)\}$ 在 $[a, A]$ 上一致收敛, 且有 $b \leq \varphi_n(x) \leq B$, 证明: $F_n(x) = f(x, \varphi_n(x))$ 在 $[a, A]$ 上一致收敛.\par
% p164-b16; source PDF page 164
\blockanchor{p164-b16}
\question{9}(10 分) 设二元函数 $f(x, y) = \varphi(|xy|)$，其中 $\varphi(0) = 0$，且在 $u = 0$ 附近满足 $|\varphi(u)| \leq u^2$，证明：$f(x, y)$ 在原点 $(0, 0)$ 处可微.\par
% p164-b17; source PDF page 164
\blockanchor{p164-b17}
\question{10}(10 分) 计算 $\iint_{D} |\sin(x - y)| \, \mathrm{d}x \, \mathrm{d}y$, 其中 $D = \{(x, y) \mid 0 \leq x \leq y \leq 2\pi\}$.\par
% p164-b18; source PDF page 164
\blockanchor{p164-b18}
\question{11}(15 分) 设 $f(x)$ 在 $[0, +\infty)$ 上连续, 且 $\int_{0}^{+\infty} f(x) \mathrm{d}x$ 收敛, 证明: 存在 $\{x_{n}\} \subset [0, +\infty)$, 满足 $\lim_{n \to \infty} x_{n} = +\infty$, 且 $\lim_{n \to \infty} f(x_{n}) = 0$.\par
% p164-b19; source PDF page 164
\blockanchor{p164-b19}
(另一个版本是: 设 $f(x)$ 在 $[0, +\infty)$ 上可导, 且 $\int_{0}^{+\infty} f(x) \mathrm{d}x$ 收敛, 证明: 存在 $\{x_n\} \subset [0, +\infty)$, 满足 $\lim_{n \to \infty} x_n = +\infty$, 且 $\lim_{n \to \infty} f'(x_n) = 0$\par
% p165-b00; source PDF page 165
\par\Needspace{6\baselineskip}
\blockanchor{p165-b00}
\question{12}(15 分) 计算曲面积分\par
\nopagebreak[4]
% p165-b01; source PDF page 165
\blockanchor{p165-b01}
\displayformula{\iint_ {S} (x - y) \mathrm{d} x \mathrm{d} y + x (y - z) \mathrm{d} y \mathrm{d} z.}
% p165-b02; source PDF page 165
\blockanchor{p165-b02}
\displayformula{\left(\text {另一个版本为} \iint_ {S} x y ^ {2}   \mathrm{d} x \mathrm{d} y + x (y - z)   \mathrm{d} y \mathrm{d} \dot {z}\right)}
% p165-b03; source PDF page 165
\blockanchor{p165-b03}
其中 S 是圆柱面  $x^{2} + y^{2} = 1$  被 z = 0, z = 3 所截部分的外侧.\par
